% Generated draft: requires mathematical and source review.
\documentclass[11pt]{article}
\usepackage[margin=25mm]{geometry}
\usepackage{fontspec}
\setmainfont{Latin Modern Roman}
\usepackage{amsmath,amssymb,mathtools,unicode-math}
\setmathfont{Latin Modern Math}
\usepackage{enumitem,needspace,xurl,hyperref}
\setlength{\parindent}{0pt}
\setlength{\parskip}{5pt}
\setlength{\emergencystretch}{3em}
\newcommand{\N}{\mathbb{N}}
\newcommand{\Z}{\mathbb{Z}}
\newcommand{\Q}{\mathbb{Q}}
\newcommand{\R}{\mathbb{R}}
\newcommand{\C}{\mathbb{C}}
\newcommand{\F}{\mathbb{F}}
\newcommand{\A}{\mathbb{A}}
\newcommand{\PP}{\mathbb{P}}
\newcommand{\E}{\mathbb{E}}
\newcommand{\eps}{\varepsilon}
\newcommand{\dd}{\,\mathrm{d}}
\begin{document}
\section*{UnsolvedMath ID: 20001282}
\textbf{Source problem code:} catalogue entry
\subsection{Definitions and mathematical statement}
For a centrally symmetric convex body translated to center $0$, let $K^\circ=\{y:\langle x,y\rangle\le1\text{ for every }x\in K\}$ and $\mathcal P(K)=\operatorname{vol}(K)\operatorname{vol}(K^\circ)$. The question attributed to G. Kuperberg asks whether, in every fixed combinatorial-type class of centrally symmetric $3$-polytopes except cube and octahedron types, $\mathcal P$ has no local minimum. It also asks whether each such class has one critical orbit modulo affine transformations, necessarily a local maximum. Locality is relative to the fixed face-lattice realization space; a precise differentiable critical-point convention is not supplied.

The attached aid treats a restricted family. A marked centrally symmetric planar hexagon has an affine normal form
\[
H_{p,q}=\operatorname{conv}\{\pm(p,q),\ \pm(0,1),\ \pm(-1,0)\},
\quad p,q>0,\quad p+q>1,\quad|p-q|<1.
\]
It gives the exact planar volume product
\[
\mathcal P(H_{p,q})=(p+q+1)\left(4-\frac{(p+q-1)^2}{pq}\right)\in(8,9].
\]
There is no local minimum in this parameter domain, and its unique critical affine orbit is the affine-regular hexagon, represented by $p=q=1$, a strict maximum modulo affine equivalence.

For $a>0$, the geometric product prism $P=H\times[-a,a]$ satisfies
\[
\mathcal P(P)=\frac43\mathcal P(H).
\]
The aid therefore supplies a same-face-lattice lowering curve at every centrally symmetric geometric product hexagonal prism and, by polarity, every polar product-dual hexagonal bipyramid. Such points cannot be local minima even in the full realization stratum. In the product submoduli the sharp range is $(32/3,12]$, and there is a unique critical affine orbit, a strict maximum in that submoduli.

These assertions do not show that every centrally symmetric realization of those combinatorial types is a product or product-dual, nor establish critical-orbit uniqueness or maximality in their full realization strata. They do not settle Kuperberg's assertion for every allowed combinatorial type.
\subsection{Short English statement}
Study local volume-product extrema within a fixed centrally symmetric polytope type. The supplied aid rules out local minima at product hexagonal prisms and their polar duals, but confines its unique-maximum statement to the product subfamily.
\subsection{Sources}
\begin{itemize}
\item Local UnsolvedMath registry, permanent ID $20001282$, source code catalogue entry; canonical \texttt{aim-convex-geometry-notes.json}, \texttt{notes[13]}.
\item American Institute of Mathematics, ``Mahler's conjecture and duality in convex geometry,'' question attributed to G. Kuperberg: \url{https://aimath.org/WWN/mahlerduality/mahlerduality.pdf}. The metadata calls it source item 8, while verification says a leading digit was lost and the terminal 1 is page spill; the corrected full number is not reproduced here. No PDF inspection is claimed, and original literature and remarks fields are empty.
\item Attached research attempt 1: partial result, recorded run status open, rigor proved, novelty confidence low; registry status partially solved, classification PARTIAL-PROGRESS. Its explicit hexagon/prism results are machine generated and not peer reviewed; independently verify before citation. No authoritative theorem bibliography is supplied.
\end{itemize}
\end{document}
